% Propietario conservado: /Users/ruben/Documents/New project/output/EXTREMA_MEDIA_RAZON_RECIPROCIDAD_ES_EN_20260918/ARTICULO/ES/source/libro/05_levantamiento.tex
\chapter{Levantamiento unitario del refinamiento y transporte de lectores}
\label{x_reciprocidad:lib05:levantamiento}

La biyección por residuos de la \cref{x_reciprocidad:lib03:teo-residuo} y el
balance bilateral de la \cref{x_reciprocidad:lib04:teo-balance} admiten una
composición explícita. Una terna aritmética admisible, que contiene
una pareja y su acarreo, determina una pareja refinada. La
linealización de esta biyección sobre bases ortonormales proporciona
un transporte unitario; dos transportes de ese tipo pueden actuar
como las dos realizaciones del balance bilateral.

La realización se establece después de los objetos generados por
APP--TRIT--TPK. Su dominio combinatorio comprende las parejas y los
acarreos que satisfacen las condiciones escritas abajo. La norma
conservada por su linealización es la suma de cuadrados de amplitudes
sobre esas etiquetas; la conservación aritmética de
\(L_c(x,y)=(Bx-c,By+c)\) continúa siendo la suma de las dos
coordenadas tras reescalar. La composición relaciona ambas
conservaciones mediante una aplicación concreta.

\section{Biyección de los estados aritméticos admisibles}
\label{x_reciprocidad:lib05:biyeccion}

Sean \(B\geq2\) y \(M\geq1\) enteros. Definimos
\begin{align}
 \mathcal D_{B,M}
 &=\bigl\{(x,y,c)\in\mathbb Z_{\geq0}^2
                 \times\{0,\ldots,B-1\}:x+y=M,\ c\leq Bx\bigr\},
 \label{x_reciprocidad:lib05:dominio}\\
 \mathcal E_{BM}
 &=\bigl\{(X,Y)\in\mathbb Z_{\geq0}^2:X+Y=BM\bigr\}.
 \label{x_reciprocidad:lib05:codominio}
\end{align}
La condición \(c\leq Bx\) mantiene la no negatividad de la
primera coordenada de la salida. Para \(x=0\) sólo admite
\(c=0\); para \(x\geq1\) admite todo el alfabeto canónico.

\begin{theorem}[Correspondencia exacta de etiquetas]
\label{x_reciprocidad:lib05:teo-biyeccion}
La aplicación
\begin{equation}
 \ell_{B,M}:\mathcal D_{B,M}\longrightarrow\mathcal E_{BM},
 \qquad (x,y,c)\longmapsto(Bx-c,By+c)
 \label{x_reciprocidad:lib05:mapa-etiquetas}
\end{equation}
es biyectiva. Ambos conjuntos tienen cardinal \(BM+1\). Su inversa
está dada por
\begin{equation}
 c=Y\bmod B,\qquad y=\frac{Y-c}{B},\qquad x=\frac{X+c}{B}.
 \label{x_reciprocidad:lib05:inversa-etiquetas}
\end{equation}
\end{theorem}

\begin{proof}
La condición de dominio y la identidad
\(Bx-c+By+c=BM\) aseguran que la imagen pertenece a
\(\mathcal E_{BM}\). Dada una pareja de ese conjunto, el residuo
\(c\) es el único elemento del alfabeto que satisface
\(Y\equiv c\pmod B\). Como \(X+Y\) es divisible por \(B\),
\(X+c\) también lo es. La división euclídea de \(Y\) da
\(y\geq0\), y \(X+c\geq0\) da \(x\geq0\). Si \(c>0\),
\(x=0\) exigiría \(X=-c<0\), por lo que se conserva la
admisibilidad. La suma de las preimágenes es \(M\), y la
sustitución verifica la imagen original. La unicidad del residuo y
de los cocientes demuestra la biyección.

Hay \(M+1\) etiquetas con \(c=0\) y \(M\) con cada acarreo
positivo; por tanto el cardinal es
\((M+1)+(B-1)M=BM+1\), coincidente con el de la salida.
\end{proof}

La prueba explicita el extremo de frontera del dominio, condición
necesaria para la igualdad de cardinales. La nueva etiqueta de acarreo
es una parte tipada de la entrada al refinamiento.

\section{Unitariedad sobre las amplitudes de las etiquetas}
\label{x_reciprocidad:lib05:unitariedad}

Tomemos los espacios de Hilbert reales o complejos con medida de conteo
\begin{equation}
 \mathscr H_D=\ell^2(\mathcal D_{B,M}),\qquad
 \mathscr H_E=\ell^2(\mathcal E_{BM}),
 \label{x_reciprocidad:lib05:hilbert}
\end{equation}
y sus bases ortonormales etiquetadas. Definimos linealmente
\begin{equation}
 J_{B,M}|x,y,c\rangle=|Bx-c,By+c\rangle.
 \label{x_reciprocidad:lib05:j}
\end{equation}
Se escribirá \(J\) cuando los parámetros estén fijados.

\begin{proposition}[Levantamiento unitario]
\label{x_reciprocidad:lib05:prop-unitaria}
El operador \(J:\mathscr H_D\to\mathscr H_E\) es unitario:
\(J^*J=I_D\), \(JJ^*=I_E\). El adjunto recupera cada etiqueta
mediante \eqref{x_reciprocidad:lib05:inversa-etiquetas}.
\end{proposition}

\begin{proof}
La \cref{x_reciprocidad:lib05:teo-biyeccion} asegura que \(J\) transforma
biyectivamente dos bases ortonormales completas. Para
\(\psi=\sum_{d\in\mathcal D}\psi_d|d\rangle\),
\[
 \|J\psi\|^2
 =\sum_{e\in\mathcal E}|\psi_{\ell_{B,M}^{-1}(e)}|^2
 =\sum_{d\in\mathcal D}|\psi_d|^2=\|\psi\|^2.
\]
La sobreyectividad prueba las dos identidades unitarias y la
biyección inversa describe el adjunto sobre cada vector de base.
\end{proof}

La etiqueta \(|73,27,1\rangle\) tiene norma uno; los enteros
\(73,27,1\) describen sus valores aritméticos, no sus amplitudes.
La construcción representa un procesamiento reversible de etiquetas
o amplitudes. Una realización física concreta conserva, además,
el protocolo que prepare y transporte tales amplitudes.

\section{Lectores diagonales de proporción y acarreo}
\label{x_reciprocidad:lib05:diagonales}

Para una función real \(h\) sobre un conjunto finito, denotemos
por \(M_h\) su operador diagonal: cada vector de base se multiplica
por el valor de \(h\) en su etiqueta. La biyección satisface
\begin{equation}
 J^*M_hJ=M_{h\circ\ell_{B,M}}.
 \label{x_reciprocidad:lib05:lector-diagonal}
\end{equation}
Se demuestra actuando sobre todos los vectores de la base.

Definamos en la salida
\begin{equation}
 f_E(X,Y)=\frac{X}{BM},\quad
 g_E(X,Y)=\frac{Y}{BM},\quad
 \gamma_E(X,Y)=Y\bmod B,
 \label{x_reciprocidad:lib05:lectores-salida}
\end{equation}
y en la entrada \(f_D=x/M\), \(g_D=y/M\), \(\gamma_D=c\).
La identidad \eqref{x_reciprocidad:lib05:lector-diagonal} proporciona
\begin{align}
 J^*M_{f_E}J&=M_{\,x/M-c/(BM)},\nonumber\\
 J^*M_{g_E}J&=M_{\,y/M+c/(BM)},\nonumber\\
 J^*M_{\gamma_E}J&=M_c.
 \label{x_reciprocidad:lib05:transporte-fracciones}
\end{align}
Los dos primeros lectores transportan la transferencia de unidad.
Como \(f_E+g_E=1\), su suma es \(I_E\) y su compresión es
\(I_D\). El tercer lector conserva exactamente el acarreo como
observable entero, con su representación residual en la salida.

\section{Rutas completas a profundidad arbitraria}
\label{x_reciprocidad:lib05:rutas}

Sea \(\mathcal P_n(B,M)\) el conjunto de rutas de longitud \(n\)
que parten de todas las parejas no negativas de suma \(M\) y, en
cada etapa, satisfacen las condiciones de
\eqref{x_reciprocidad:lib05:dominio} para la suma correspondiente. Una ruta conserva
su pareja inicial y sus \(n\) acarreos admisibles. Sea
\(\Phi_n\) la aplicación que asigna a cada ruta su pareja final.

\begin{theorem}[Recuperación de rutas a toda profundidad finita]
\label{x_reciprocidad:lib05:teo-rutas}
Para todo \(n\geq0\),
\begin{equation}
 \Phi_n:\mathcal P_n(B,M)\longrightarrow\mathcal E_{B^nM}
 \quad\hbox{es biyectiva},\qquad
 |\mathcal P_n(B,M)|=B^nM+1.
 \label{x_reciprocidad:lib05:biyeccion-rutas}
\end{equation}
Su linealización
\begin{equation}
 J_n:\ell^2(\mathcal P_n(B,M))\longrightarrow
          \ell^2(\mathcal E_{B^nM}),\qquad
 J_n|\pi\rangle=|\Phi_n(\pi)\rangle
 \label{x_reciprocidad:lib05:jn}
\end{equation}
es unitaria.
\end{theorem}

\begin{proof}
Desde una pareja de suma \(B^nM\), el residuo de
\eqref{x_reciprocidad:lib05:inversa-etiquetas} recupera el último acarreo y una
única pareja previa de suma \(B^{n-1}M\). Repetir \(n\) veces
produce una única ruta cuya entrada suma \(M\). La composición
directa recupera el extremo original, lo que prueba biyección e
inversa. La cuenta también se obtiene desde las entradas: las
\(M\) parejas con \(x_0\geq1\) admiten \(B^n\) palabras cada
una, mientras que \((0,M)\) sólo admite la palabra de todos
ceros. El total es \(MB^n+1\). La transformación biyectiva de
bases ortonormales demuestra la unitariedad de \(J_n\).
\end{proof}

Conocida la profundidad, los lectores sobre la pareja final son
\begin{equation}
 c_j=\left\lfloor\frac{Y}{B^{n-j}}\right\rfloor\bmod B,
 \quad y_0=\left\lfloor\frac{Y}{B^n}\right\rfloor,
 \quad x_0=M-y_0,
 \qquad 1\leq j\leq n.
 \label{x_reciprocidad:lib05:lectores-ruta}
\end{equation}
Estas funciones son diagonales y recuperan la memoria completa de
las operaciones finitas. La compatibilidad con prolongación es
\begin{equation}
 \Phi_{n+1}(\pi,c)=L_c(\Phi_n(\pi)).
 \label{x_reciprocidad:lib05:compatibilidad-prolongacion}
\end{equation}
Al pasar de \(n\) a \(n+1\), se incorpora la nueva etiqueta
admisible de acarreo. Los espacios aumentan su dimensión y la
etiqueta de frontera sólo admite \(c=0\); por ello el dominio del
paso siguiente conserva la admisibilidad, en lugar de reemplazarla
por un producto irrestricto con \(B\) dígitos.

\section{Entrada fijada e imagen exacta}
\label{x_reciprocidad:lib05:entrada-fija}

Si se fija una pareja inicial \((x_0,y_0)\), con \(x_0\geq1\),
hay \(B^n\) rutas y su imagen es exactamente
\begin{equation}
 \bigl\{(B^nx_0-C,B^ny_0+C):0\leq C\leq B^n-1\bigr\}.
 \label{x_reciprocidad:lib05:imagen-fija}
\end{equation}
La \cref{x_reciprocidad:lib03:cor-palabra} demuestra que los \(B^n\)
acumulados son distintos y recuperan las respectivas palabras. La
linealización es unitaria hacia el subespacio generado por esta
imagen. Si \(x_0=0\), la imagen consta de una sola etiqueta.

Los cilindros de la \cref{x_reciprocidad:lib03:teo-cilindros} corresponden a una
entrada fijada. El dominio completo de
\eqref{x_reciprocidad:lib05:biyeccion-rutas} reúne todas las parejas de suma \(M\).
Ambas construcciones poseen inversas explícitas, pero sus imágenes
y cardinales son diferentes. La profundidad arbitraria expresa la
validez de la prueba para todo \(n\); el límite de la lectura y
sus fronteras sigue siendo el establecido en la
\cref{x_reciprocidad:lib03:limite}.

\section{Acoplamiento bilateral del refinamiento}
\label{x_reciprocidad:lib05:bilateral}

Sea \(R\) una permutación previamente declarada de
\(\mathcal D_{B,M}\), linealizada unitariamente en
\(\mathscr H_D\). Definimos los dos transportes
\begin{equation}
 \mathcal B=J,\qquad \mathcal C=JR:
            \mathscr H_D\longrightarrow\mathscr H_E.
 \label{x_reciprocidad:lib05:dos-transportes}
\end{equation}
La permutación forma parte de la especificación del procedimiento
y se fija antes de evaluar una lectura de salida.
Para \(a>0\), \(p=a^2+a+1\), \(N=(2a+1)p\),
\begin{equation}
 Q_-=J(aI-R),\qquad Q_+=J((a+1)I+R).
 \label{x_reciprocidad:lib05:q-refinamiento}
\end{equation}

\begin{proposition}[Distribución y recuperación de amplitudes refinadas]
\label{x_reciprocidad:lib05:prop-bilateral}
Los operadores \eqref{x_reciprocidad:lib05:q-refinamiento} cumplen
\begin{equation}
 (a+1)Q_-^*Q_-+aQ_+^*Q_+=NI_D,
 \label{x_reciprocidad:lib05:balance}
\end{equation}
y la aplicación
\begin{equation}
 W_a\psi=\frac1{\sqrt N}
    \bigl(\sqrt{a+1}\,Q_-\psi,\sqrt a\,Q_+\psi\bigr)
 \label{x_reciprocidad:lib05:w}
\end{equation}
es isométrica. Para los canales sin normalizar
\(u=Q_-\psi\), \(v=Q_+\psi\),
\begin{equation}
 J\psi=\frac{u+v}{2a+1},\qquad
 \psi=J^*\frac{u+v}{2a+1}.
 \label{x_reciprocidad:lib05:inversa-bilateral}
\end{equation}
\end{proposition}

\begin{proof}
Las dos aplicaciones \(\mathcal B,\mathcal C\) son unitarias.
Los términos cruzados de los Grams de \(Q_-\) y \(Q_+\) tienen
coeficientes \(-a\) y \(a+1\), respectivamente. Las ponderaciones
del balance los cancelan, mientras los términos diagonales suman
\(N\), como en la \cref{x_reciprocidad:lib04:teo-balance}. La normalización
demuestra la isometría. La suma de los canales es
\((2a+1)J\psi\), y la unitariedad de \(J\) da la inversa.
\end{proof}

Desde los canales normalizados se recupera directamente la entrada
por \(W_a^*\), con norma inversora uno; este adjunto ya incorpora
\(J^*\). La fórmula alternativa
\eqref{x_reciprocidad:lib05:inversa-bilateral} separa la recuperación de \(J\psi\)
por suma de la aplicación de \(J^*\), cuya acción sobre etiquetas
se calcula mediante residuos. Quedan compuestos el refinamiento
aritmético reversible y el transporte bilateral de amplitudes.

\begin{proposition}[Compatibilidad en el espacio refinado]
\label{x_reciprocidad:lib05:prop-compatibilidad}
Un par sin normalizar \((u,v)\) pertenece a la imagen de
\(\psi\mapsto(Q_-\psi,Q_+\psi)\) si y sólo si
\begin{equation}
 \bigl[(a+1)I_E+JRJ^*\bigr]u
 =\bigl[aI_E-JRJ^*\bigr]v.
 \label{x_reciprocidad:lib05:compatibilidad}
\end{equation}
\end{proposition}

\begin{proof}
La igualdad equivale a
\(JRJ^*(u+v)=av-(a+1)u\). Si se cumple, definamos
\(\psi=J^*(u+v)/(2a+1)\). Al sustituir en
\eqref{x_reciprocidad:lib05:q-refinamiento}, esa identidad da
\(Q_-\psi=u\) y \(Q_+\psi=v\). La necesidad resulta de
insertar los canales verdaderos. Todas las composiciones actúan
en \(\mathscr H_E\), por lo que sus tipos están definidos.
\end{proof}

\section{Ley de los observables sobre los dos canales}
\label{x_reciprocidad:lib05:ley-lectores}

\begin{theorem}[Compresión bilateral de un lector refinado]
\label{x_reciprocidad:lib05:teo-lector}
Para \(F=F^*\) sobre \(\mathscr H_E\), escribamos
\(F_J=J^*FJ\). Entonces
\begin{equation}
 W_a^*\operatorname{diag}(F,F)W_a
 =\frac{a(a+1)F_J+R^*F_JR}{p}.
 \label{x_reciprocidad:lib05:lector-bilateral}
\end{equation}
El lector \(F_J\) se conserva exactamente si y sólo si conmuta
con \(R\).
\end{theorem}

\begin{proof}
Expandir \((a+1)Q_-^*FQ_-+aQ_+^*FQ_+\) cancela los
términos mixtos. El coeficiente de \(\mathcal B^*F\mathcal B\)
es \(a(a+1)(2a+1)\); el de
\(\mathcal C^*F\mathcal C\) es \(2a+1\). Dividir por
\(N\), sustituir \(\mathcal C=JR\) y usar
\(F_J=J^*FJ\) prueba la fórmula. La igualdad con \(F_J\)
equivale a \(R^*F_JR=F_J\), o a \([F_J,R]=0\) por
unitariedad.
\end{proof}

En el caso \(a=8\), los pesos son \(72/73\) y \(1/73\).
El lector identidad recupera la norma. Un lector no constante
puede cambiar aun cuando la norma se conserve; su variación está
calculada por \eqref{x_reciprocidad:lib05:lector-bilateral}.
El operador que conserva exactamente un lector original \(A\)
en la imagen completa es \(W_aAW_a^*\), cuya identidad de imagen
es \(W_aW_a^*\). Leer \(\operatorname{diag}(F,F)\) corresponde
a la compresión específica del teorema, y no sustituye
automáticamente ese observable transportado.

\section{Ejemplo aritmético de dimensión mil uno}
\label{x_reciprocidad:lib05:ejemplo}

Tomemos \(B=10\), \(M=100\). Los espacios de etiquetas tienen
dimensión \(1001\), y
\begin{equation}
 d_1=(73,27,1)\longmapsto e_1=(729,271).
 \label{x_reciprocidad:lib05:ejemplo-etiqueta}
\end{equation}
Los lectores diagonales dan
\begin{equation}
 f_E(e_1)=\frac{729}{1000}
 =\frac{73}{100}-\frac1{1000},\qquad
 g_E(e_1)=\frac{271}{1000},\qquad \gamma_E(e_1)=1.
 \label{x_reciprocidad:lib05:ejemplo-lectores}
\end{equation}
Declaremos \(R\) como la transposición de las etiquetas
admisibles \(d_1=(73,27,1)\) y \(d_2=(73,27,2)\), dejando
fijas las demás. La segunda imagen es \(e_2=(728,272)\).
La regla de transposición se ha fijado sobre las etiquetas, antes
de cualquier comparación con una constante.

Para \(a=8\) y \(\psi=|d_1\rangle\),
\begin{equation}
 Q_-\psi=8|e_1\rangle-|e_2\rangle,\qquad
 Q_+\psi=9|e_1\rangle+|e_2\rangle.
 \label{x_reciprocidad:lib05:ejemplo-canales}
\end{equation}
Las normas cuadradas son \(65\) y \(82\), y
\begin{equation}
 9\cdot65+8\cdot82=585+656=1241=17\cdot73.
 \label{x_reciprocidad:lib05:ejemplo-balance}
\end{equation}
La suma de los canales es \(17|e_1\rangle\). La inversa
\eqref{x_reciprocidad:lib05:inversa-bilateral} y el residuo recuperan
\(|d_1\rangle\), incluido el acarreo uno.

Si se evalúa \(F=M_{f_E}\) diagonalmente en ambos canales
normalizados, la \cref{x_reciprocidad:lib05:teo-lector} da
\begin{equation}
 \begin{split}
 \langle W_8\psi,\operatorname{diag}(F,F)W_8\psi\rangle
 &=\frac{72\cdot729+728}{73\cdot1000}\\
 &=\frac{729}{1000}-\frac1{73000}.
 \end{split}
 \label{x_reciprocidad:lib05:ejemplo-fraccion}
\end{equation}
La lectura complementaria es uno menos esta cantidad. Para el
lector diagonal de acarreo se obtiene
\begin{equation}
 \frac{72\cdot1+2}{73}=\frac{74}{73}.
 \label{x_reciprocidad:lib05:ejemplo-acarreo}
\end{equation}
El acarreo original sigue siendo recuperable exactamente: el
observable \(W_8M_cW_8^*\) tiene sobre \(W_8|d_1\rangle\)
el valor propio uno. El valor \(74/73\) es la respuesta del
lector diagonal aplicado a ambos canales, cuyo transporte relativo
está especificado por \eqref{x_reciprocidad:lib05:lector-bilateral}. El ejemplo
distingue conservación de la información y variación de una lectura
particular.

\section{Articulación con el registro generado y dominios de aplicación}
\label{x_reciprocidad:lib05:dominios}

La unitaria \(J\) actúa sobre \(BM+1\) etiquetas, mientras que
la realización de \(K\) tiene doce coordenadas. En el ejemplo,
\(1001\) cuenta parejas refinadas de suma mil, y doce cuenta
ventanas del registro producido. Ambas dimensiones conservan su
origen combinatorio específico.

Una composición particular entre esos dos portadores puede formularse
mediante un selector material de doce rutas distintas, con su
procedencia y sus operadores. Si ese selector proporciona una
isometría
\(J_{12}:\mathbb R^{12}\to\ell^2(\mathcal D_{B,M})\), entonces
\(JJ_{12}\) es isométrica, y el entrelazamiento de los lectores se
comprueba sobre ella. Esta es la hipótesis de esa composición
adicional, no una identificación implícita por el número de
coordenadas. Las aplicaciones bilaterales ya definidas directamente
sobre \(K\) en la \cref{x_reciprocidad:lib04:bilateral} conservan sus dominios
completos sin utilizar ese selector adicional.

El resultado de esta sección reúne una biyección aritmética, su
unitaria, lectores diagonales entrelazados, recuperación de rutas a
cualquier profundidad finita y una composición bilateral explícita.
El dominio de rutas es el definido en
\eqref{x_reciprocidad:lib05:biyeccion-rutas}. Las aplicaciones a historias
filtradas por TPK conservan sus reglas de supervivencia y de
prolongación; la correspondencia de etiquetas prueba aquí el
transporte declarado, sin reemplazar esos operadores por un dominio
aritmético irrestricto.
